A natural application of neuron addition is to non-stationary data distributions, such as Continual Learning. Instead of a single task, represented by its dataset $\mathcal{D}$, we have a sequence of tasks $\mathcal{D}_1, \dots, \mathcal{D}_T$, and want a model that performs well on all tasks $\{\mathcal{D}_1, \dots, \mathcal{D}_T\}$. In addition to the objective of maximising performance on the current task $\mathcal{D}_i$, the goal is also to prevent relapse on previous tasks, known as \emph{catastrophic forgetting}~\citep{lange_continual_2021}. The trade-off between performing well on the new task while retaining performance on the old tasks is known as the \emph{plasticity-stability dilemma}. Parameter isolation methods, which dedicate specific parameters to individual tasks, are a particularly strong family of Continual Learning algorithms, and Neural Network growth is naturally adapted to this, as in Progressive Neural Networks~\citep{rusu_progressive_2016}, which add new layers for each task.

\textbf{Continual Learning.} By aligning the architectures and weights $(A_t, \theta_t)$ with the current task $\D_t$, growing allows one to reuse existing weights and only add parameters for the next task $\D_{t+1}$ when necessary. Some methods enforce sparsity in the weights ~\citep{yoon_lifelong_2018,yang_grown_2021} or temporary pruning ~\citep{hung_CPG_2019,wu_firefly_2020} --- disabling the weights to be able to adapt them when training the next task without introducing semantic drift. This results in supernets where each task uses a subset of the weights. Based on ablation studies~\citep{yoon_lifelong_2018}, growth avoids catastrophic forgetting and adds flexibility to the parameters while maintaining smaller models. The \textbf{Learn-to-Grow} framework~\citep{li_learn_2019} uses NAS to balance performance and accuracy by reusing existing weights as much as possible and replacing them when necessary.

\textbf{Reinforcement Learning.} Loss of plasticity~\citep{dohare_loss_2024} is also a significant problem in Reinforcement Learning: the value function and/or policy are represented by Neural Networks and quickly overfit to initial observations, failing to adapt to new data as training progresses. Proposed solutions include periodically resetting some or all of the network~\citep{nikishin_primacy_2022,sokar_dormant_2023}, or using various forms of regularization~\citep{bhatt_crossq_2024,galliciSimplifyingDeepTemporal2025}. 
Recently, growth methods have been used to maintain plasticity with a fixed growth schedule~\citep{kangForgetandGrowStrategyDeep2025} or with sparse grow-prune methods~\citep{liu_neuroplastic_2025}.
